\medskip

\begin{tabularx}{\linewidth}{|X|}\hline
\textbf{Principaux domaines abordés:} Suites ; Fonctions, Fonction logarithme.\\ \hline
\end{tabularx}

\medskip

\emph{Cet exercice est un questionnaire à choix multiples.\\
Pour chacune des questions suivantes, une seule des quatre réponses proposées
est exacte.\\
Une réponse exacte rapporte un point. Une réponse fausse, une réponse multiple ou
l'absence de réponse à une question ne rapporte ni n'enlève de point.\\
Pour répondre, indiquer sur la copie le numéro de la question et la lettre de la réponse choisie. Aucune justification n'est demandée.}

\bigskip
\begin{enumerate}
\item Un récipient contenant initialement 1 litre d'eau est laissé au soleil.

Toutes les heures, le volume d'eau diminue de 15\,\%.

Au bout de quel nombre entier d'heures le volume d'eau devient-il inférieur à un quart de litre ?

\begin{center}
\begin{tabularx}{\linewidth}{*{4}{X}}
\textbf{a.~~} 2 heures &\textbf{b.~~} 8 heures .&\textbf{c.~~} 9 heures
&\textbf{d.~~} 13 heures
\end{tabularx}
\end{center}
Enlever 15\,\%, revient à multiplier par 0,85 $\left(1 - \dfrac{15}{100} = \dfrac{85}{100} = 0,85 \right)$.

Il faut donc trouver $n \in \N$ tel que $1 \times 0,85^n < 0,25$ soit en prenant le logarithme népérien $n\ln 0,95 < \ln 0,25 \iff n > \dfrac{\ln 0,25}{\ln 0,85}$.

Or $\dfrac{\ln 0,25}{\ln 0,85} \approx 8,5$ : il faut donc attendre au moins la 9\up{e} heure.
\item  On considère la fonction $f$ définie sur l'intervalle $]0~;~+\infty[$ par $f(x) = 4\ln (3x)$.

Pour tout réel $x$ de l'intervalle $]0~;~+\infty[$ , on a :

\begin{center}
\begin{tabularx}{\linewidth}{*{2}{X}}
\textbf{a.~~} $f(2x) = f(x) + \ln (24)$&\textbf{b.~~} $f(2x) = f(x) + \ln (16)$\\
\textbf{c.~~} $f(2x) = \ln (2) + f(x)$& \textbf{d.~~} $f(2x) = 2f(x)$
\end{tabularx}
\end{center}

$\bullet~~$ On a $f(x) + \ln (16) = 4 \ln (3x) + \ln 2^4= \ln (3x)^4  + \ln 2^4 = \ln \left[2^4 \times (3x)^4\right] = \ln \left[(6x)^4\right] = 4 \ln (6x) = f(2x)$ Réponse \textbf{b.}

$\bullet~~$ Autre solution : $f(2x) = 4\ln (3 \times 2x) = 4 \ln (2 \times 3x) = 4\ln 2 + 4\ln (3x) = \ln 2^4 + 4\ln (3x) =$

$ \ln 16 + 4\ln (3x) = f(x) + \ln 16$.
\item On considère la fonction $g$ définie sur l'intervalle $]1~;~+\infty[$ par : 

\[g(x) = \dfrac{\ln (x)}{x - 1}.\]


On note $\mathcal{C}_g$ la courbe représentative de la fonction $g$ dans un repère orthogonal. La courbe $\mathcal{C}_g$ admet :

$\bullet~~$Limite de la fonction $g$ au voisinage de plus l'infini :

$g(x) = \dfrac{\ln x}{x} \times \dfrac{x}{x - 1}$.

Or on sait que $\displaystyle\lim_{x \to +\infty} \dfrac{\ln x}{x} = 0$ et que  $\displaystyle\lim_{x \to +\infty} \dfrac{x}{x - 1} = 1$, donc par produit de limites : $\displaystyle\lim_{x \to +\infty} g(x) = 0$ : l'axe des abscisses est asymptote horizontale au voisinage de plus l'infini.

$\bullet~~$Limite de la fonction $g$ au voisinage de 1 : on a :

$g(x) = \dfrac{\ln (x)}{x - 1} = \dfrac{\ln (x) - \ln (1)}{x - 1}$, donc $\displaystyle\lim_{x \to 1} \dfrac{\ln (x) - \ln (1)}{x - 1} = [\ln (x)]'_{x = 1}$, nombre dérivé de la fonction logarithme népérien en $x = 1$, soit $\dfrac{1}{1} = 1$ : il n'y a as d'asymptote verticale au voisinage de 1. Réponse : \textbf{c.}










\end{enumerate}

Dans la suite de l'exercice, on considère la fonction $h$ définie sur l'intervalle ]0~;~2] par:

\[h(x) = x^2(1 + 2\ln (x)).\]

On note $\mathcal{C}_h$ la courbe représentative de $h$ dans un repère du plan. 

On admet que $h$ est deux fois dérivable sur l'intervalle ]0~;~2].

On note $h'$ sa dérivée et $h''$ sa dérivée seconde.

On admet que, pour tout réel $x$ de l'intervalle ]0~;~2], on a :

\[h'(x) = 4x(1 + \ln (x)).\]

\begin{enumerate}[resume]
\item Sur l'intervalle ]0~;~2], la fonction $h$ s'annule :

\begin{center}
\begin{tabularx}{\linewidth}{*{2}{X}}
\textbf{a.~~} exactement 0 fois. &\textbf{b.~~} exactement 1 fois.\\
\textbf{c.~~} exactement 2 fois. &\textbf{d.~~} exactement 3 fois.
\end{tabularx}
\end{center}
On a $h(x) = 0 \iff x^2(1 + 2\ln (x)) = 0 \iff \left\{\begin{array}{l c l}
x^2 &=&0\\
1 + 2\ln (x)&=&0
\end{array}\right. \iff \left\{\begin{array}{l c l}
x &=&0\\
1 + 2\ln (x)&=&0
\end{array}\right. \iff \left\{\begin{array}{l c l}
x &=&0\\
2\ln (x)&=&- 1
\end{array}\right. \iff \left\{\begin{array}{l c l}
x &=&0\\
\ln (x)&=&- \dfrac{1}{2}\end{array}\right.\iff \left\{\begin{array}{l c l}
x &=&0\\
x&=&\text{e}^{- \frac{1}{2}}\end{array}\right.$

Comme 0 ne peut être solution et que $\text{e}^{- \frac{1}{2}} \in ]0~;~2]$, l'équation a une solution : réponse \textbf{b.}
\item Une équation de la tangente à $\mathcal{C}_h$ au point d'abscisse $\sqrt{\text{e}}$ est:

\begin{center}
\begin{tabularx}{\linewidth}{*{2}{X}}
\textbf{a.~~} $y = \left(6\text{e}^{\frac12}\right).{} x $&
\textbf{b.~~} $y = \left(6\sqrt{\text{e}}\right).{} x + 2\text{e}$\\
\textbf{c.~~} $y = 6\text{e}^{\frac{x}{2}}$
&\textbf{d.~~} $y = \left(6\text{e}^{\frac12}\right) .{} x - 4\text{e}$.
\end{tabularx}
\end{center}
Une équation de la tangente est $y - h\left(\sqrt{\text{e}}\right) = f'\left(\sqrt{\text{e}}\right)\left(x - \sqrt{\text{e}}\right)$.

$\bullet~~$$h'\left(\sqrt{\text{e}}\right)  = 4 \times \sqrt{\text{e}}\left(1 + \ln \left(\sqrt{\text{e}}\right)\right) = 4 \times \sqrt{\text{e}} \times \left(1 + \dfrac12\right) = 6\sqrt{\text{e}}$.

$\bullet~~$$h\left(\sqrt{\text{e}}\right) = \left(\sqrt{\text{e}}\right)^2\left(1 + 2 \ln \left(\sqrt{\text{e}}\right)\right) = \text{e}\ln \left(1 +  2 \times \dfrac12\ln \text{e}\right) = \text{e}\times (1 + 1) = 2\text{e}$.

L'équation de la tangente s'écrit donc :

$y - 2\text{e} = 6\sqrt{\text{e}}\left(x - \sqrt{\text{e}}\right) \iff y = 2\text{e} + 6x\sqrt{\text{e}} - 6\text{e}\iff y = 6x\sqrt{\text{e}} - 4\text{e}$. Réponse \textbf{d.}

\item Sur l'intervalle $]0~;~2]$, le nombre de points d'inflexion de la courbe $\mathcal{C}_h$ est égal à :

\begin{center}
\begin{tabularx}{\linewidth}{*{4}{X}}
\textbf{a.~~} 0&\textbf{b.~~} 1 &\textbf{c.~~} 2 &\textbf{d.~~} 3
\end{tabularx}
\end{center}

La fonction $h'$ produit de fonctions dérivables sur ]0~;~2] est dérivable et sur cet intervalle :

$f''(x) = 4(1 + \ln (x)) + 4x \times \dfrac{1}{x} = 4 + 4\ln (x)) + 4 = 8 + 4\ln (x) = 4(2 + \ln (x)$.

De même $h'''(x) = \dfrac{4}{x}$. Sur l'intervalle ]0~;~2], $h'''(x) > 0$, donc $h''(x)$ est croissante et s'annule si $h''(x) = 0 \iff 2 + \ln (x) = 0 \iff \ln (x) = - 2 \iff x = \text{e}^{-2} \approx 0,135 \in ]0~;~2]$.

La dérivée seconde s'annule donc une seule fois sur l'intervalle ]0~;~2] en changeant de signe. Donc un seul point d'inflexion. Réponse \textbf{b.}
\end{enumerate}

\bigskip
